\section{Reproducibility and research decisions}\label{sec:verification}

\subsection{What the computations establish}

The package contains three standalone Python programs. They use only
the Python standard library and make no network requests. Their outputs
are already included as JSON and CSV. The scripts certify finite
identities and examples; the general theorems are justified by the
proofs in this report.

\begin{table}[htbp]
\centering\small
\begin{tabularx}{\textwidth}{@{}p{0.27\textwidth}YY@{}}
\toprule
Program & Checks performed & Scope\\
\midrule
\path{check_density_certificate.py} &
All 2701 midpoint tests for the 74-point set; exact rational weight and
tail inequalities; sphere cardinality and range; integer irreducible
counts and variance cutoffs. &
Certifies the finite inputs, including \(R_2^*+13<2^{476}\).
The infinite-set proof is in Sections~\ref{sec:density}--\ref{sec:refinement}.\\
\addlinespace
\path{verify_rank_moments.py} &
Enumerates \(\GL_4(\F_2)\) and its \(\Sp_4(\F_2)\) subgroup;
tests uniform and nonuniform positive-type laws, frame coefficients,
moments, and tails with exact arithmetic. &
1277 checks pass, including 636 frame-coefficient checks.
Large-\(\GL\) examples use exact orbit counts, not enumeration.\\
\addlinespace
\path{conductor23_benchmark.py} &
Compares the eta product with the cubic-splitting Euler product
through \(n=2000\), plus prime-power selector coefficients. &
303 primes and 3636 selector checks. Does not test or prove zero
locations.\\
\bottomrule
\end{tabularx}
\caption{Reproducible finite verification.}
\end{table}

The prior density checkers were also rerun. In addition to the sphere
and weighted certificates, the inherited finite binary-polynomial
confinement test examined 64,770 triples; all passed. That test does not
replace the confinement proof for arbitrary degrees.

\subsection{Formalization boundary and proposed Lean decomposition}

The original repository is a Lean development. The additional results
in this report have not been formalized or compiled in Lean during this
session. The package includes a formalization plan rather than a file
with placeholders presented as a completed proof.

A useful decomposition is the following.

\begin{enumerate}
\item \textbf{Finite witness certificates.} Define the 74-element finset,
the absence of nonconstant midpoints, and the two integer inequalities
obtained after clearing denominators. Use kernel-checked arithmetic for
the certificates.
\item \textbf{Replacement theorem.} Isolate unique smooth/rough
factorization, the squarefree confinement lemma, and the construction
of separated degree bands. This theorem should be independent of the
analytic number theory in \path{grh-lean}.
\item \textbf{Coefficient concentration.} Formalize the tilted
coefficient distribution, the mean and variance of a finite sum of
geometric variables, the signed-tail inequality, and the exact
ceiling bounds. Keep the probability argument separate from the
polynomial-ring application.
\item \textbf{Kernel moments.} Extract the positive-type coefficient
bound as an independent lemma. Add rectangular character
orthogonality, rank factorization, the frame count, and the
nonnegative-moment tail step.
\item \textbf{Hecke interface.} State the ideal-class logarithmic
derivative and the one-sided bound as separate interfaces. Prove the
Landau bridge before attempting the missing arithmetic construction.
The hypothesis \eqref{eq:hecke-missing-bound} must remain visibly a
hypothesis until it is actually supplied.
\end{enumerate}

This separation makes each prospective formal theorem reviewable and
prevents numerical certificates, inherited Lean theorems, and new
unformalized arguments from being confused.

\subsection{Recommended order of further work}

The polynomial counterexample is the nearest complete deliverable:
write it as a short independent paper, preserve the distinction between
the old counterexample and the new cutoff, and formalize its finite
certificates and confinement lemma. The present upper-density result
is already mathematically complete at the paper-proof level.

For a new analytic-number-theory theorem, the conductor-\(23\) function
is a focused test. The first substantial milestone is not another
zero computation. It is a construction whose valuations count the
principal ideal-class prime powers in \eqref{eq:class-selector}, with
short, nearly spanning forms and a matching lower covolume estimate.
The benchmark and analytic bridge specify exactly what success would
have to produce.

The matrix-moment result is suitable for a separate reusable lemma
library. Its value is clearest for group actions or positive-type
weights where stabilizers are accessible but the full compression law
is not. Specialized uniform \(\GL_n\) laws already admit better exact
answers; the report's example explicitly acknowledges this.

Broad claims about all automorphic GRH, simplicity of zeros, pair
correlation, central leading terms, or fixed integral Selmer groups
would need additional inputs absent from the transferred interfaces.
The determinant comparison controls a particular one-sided
prime-sum remainder. It supplies no automatic estimate for these
different observables. The correct next theorem should be stated
with its own arithmetic quantity and budget before attempting a
large formalization.

\subsection{Files and provenance}

\path{report.tex} and the files in \path{sections/} contain the complete
editable manuscript. \path{report.pdf} is the rendered version.
\path{build.sh} reproduces it using \path{latexmk}; \path{run_checks.sh}
runs the three exact verification programs.

\path{sources/prior_archives.csv} lists all 39 inspected archive
identities at the pinned ProveIt commit. The associated verification
JSON records the byte-size and Git-blob checks. Selected upstream
source snapshots are in \path{sources/lean_snapshots/}; they are
provided for source inspection, not as a complete buildable copy of the
original Lean project. Their source map, SHA-256 hashes, Apache license,
and notice are included. A package-level SHA-256 manifest makes
accidental changes to any shipped artifact detectable.

The review establishes a concrete negative answer to the cited
polynomial-density conjecture, together with an improved explicit
margin and two further transfer tools. It does not assert a proof
of a new Hecke RH instance or a general resolution of the other
arithmetic conjectures discussed here.
